\section{Conclusiones}
El desarrollo del sistema de gestión de tickets permitió comprobar que la integración de 
frameworks modernos como Laravel y React, junto con la aplicación de metodologías 
ágiles como Scrum, es una estrategia efectiva para resolver problemáticass de 
trazabilidad y eficiencia en el manejo de incidencias. La comparación con la literatura 
académica revisada confirma que estas herramientas y enfoques garantizan aplicaciones 
escalables, seguras y centradas en la experiencia del usuario. 

Entre los principales aportes del proyecto se destacan: la centralización de información, 
la reducción en tiempos de atención, la personalización de roles y la posibilidad de 
generar métricas confiables para la toma de decisiones. Asimismo, el sistema constituye 
una base sólida sobre la cual es posible implementar mejoras futuras o implementar 
nuevas funciones, como la ampliación hacia plataformas iOS mediante React Native, la 
creación de módulos de inteligencia artificial y la exploración de arquitecturas basadas en 
microservicios. 

Este artículo evidencia que el uso de frameworks modernos no es una moda pasajera, 
sino una necesidad en el desarrollo de software competitivo y orientado a la calidad. 
Igualmente, se demuestra que las metodologías ágiles, aplicadas con criterio y disciplina, 
potencian la adaptabilidad, funcionamiento y aseguran productos más acordes con las 
necesidades reales de los usuarios finales. 